\section{Security、Privacy 与 Abuse：代码和 Token 都是高价值资产}

AI coding 平台同时接触私有源代码和昂贵模型 quota。安全不仅是“数据库加密”，还包括 workspace ignore、access scope、retention、provider data policy、agent tool approval、secrets 与 audit；反滥用也不是独立风控，而是保护 inference capacity 和 unit economics 的基础设施。

\subsection{代码数据流与 shared responsibility}

\term{envelope encryption}（信封加密）通常用 data encryption key 加密数据，再用 key encryption key 保护前者；client-held 或 tenant-controlled key 可以缩小 server compromise 的可读范围，但 key lifecycle、metadata、access log 和 provider path 仍需治理。Current Cursor docs 还区分 Privacy Mode、retention 与 enterprise controls。

\lecturefigure{15-security-flow.png}{Code security 由 workspace、index、Cursor service 与 model provider 多层共同承担。}{本地字幕 41:40--43:30；Cursor 当前 privacy/security 文档。}

\begin{knowledgebox}{读图：Encryption 不能替代五件事}
Access control 决定谁能读；retention 决定保留多久；audit 记录谁做了什么；data minimization 减少不必要传输；provider governance 约束第三方处理。加密只保护特定静态/传输风险，不会阻止合法身份误用或 agent 读取 workspace secrets。
\end{knowledgebox}

\begin{warningbox}{Agent tool approval 需要风险分级}
Read-only search、local edit、terminal command、network request 和 deployment 的影响不同。默认策略应按 tool、target、workspace trust 和 user intent 分级；审批 UI 要显示将执行什么，而不是让用户盲目点击“允许”。
\end{warningbox}

\subsection{Free-token arbitrage 与资源控制}

课堂描述攻击者利用免费或补贴账户获取昂贵 frontier token。系统要把 identity、plan、device/payment signal、per-user/model quota、velocity 和 anomaly 组合，而不是只依赖 CAPTCHA。\term{abuse prevention} 的目标是限制损失并保留申诉，不是无限增加普通用户摩擦。

\lecturefigure{14-abuse-controls.png}{Free token 形成可被自动化套利的 unit-economics attack surface。}{本地字幕 36:40--38:10、43:30--45:00；概念重绘。}

\begin{knowledgebox}{读图：Quota 必须映射真实成本}
不同模型的 token price、context、cache 和 tool call 成本不同；统一 request count 会低估高成本路径。Budget 应支持 token/model 权重、burst、daily cap、organization pool 与 suspicious velocity，并在 throttling 时向合法用户解释原因。
\end{knowledgebox}

\begin{warningbox}{反滥用不可把高强度开发误判为攻击}
大型 refactor、CI agent 或团队 rollout 可能产生异常高流量。Detection 应结合 account age、payment、organization、repetition pattern、concurrency 和 model mix，并提供 temporary verification 或 enterprise quota，而不是直接永久封禁。
\end{warningbox}

\teachervoice{课堂把 security 和 abuse 放在同一场问答里很合理：私有代码泄漏损害信任，token 被套利损害经济性；两者都会让产品无法持续。安全团队和平台团队必须共享 request、identity、cost 与 incident telemetry。}

\subsection{本章小结}

AI coding 的安全边界跨越 endpoint、index、service 和 model provider。Privacy controls、agent approvals 与 abuse budget 需要共同设计，才能同时保护代码、用户控制和计算资源。

